% answers_short.tex —— 2013 年·重庆卷·理科 · 简短答案速查
% 由 tools/make_short_answers.py 自 answers.tex 的 \gksolhead/\gkanswer 提取，
% 并与 2013/2013_数学真题_重庆理_答案要点.md 逐题比对后生成，勿手改（改源 MD 后重跑）。

\begin{gknotice}
\textbf{使用说明}：本节只给最终答案便于快速核对，完整推导见解析版。答案为编写者独立推导，\textbf{非官方标准答案}。
\end{gknotice}

\gkpart{选择题}
\begin{center}\begin{tabular}{ccccccccccc}
\toprule
题号 & 1 & 2 & 3 & 4 & 5 & 6 & 7 & 8 & 9 & 10 \\
\midrule
答案 & D & D & B & C & C & A & A & B & C & D \\
\bottomrule
\end{tabular}\end{center}

\gkpart{填空题}
\begin{itemize}[leftmargin=2.6em]
  \item[\textbf{11.}] $\sqrt5$。$|z|=\dfrac{|5\mathrm i|}{|1+2\mathrm i|}=\dfrac5{\sqrt5}=\sqrt5$。
  \item[\textbf{12.}] $64$。由 $a_2^2=a_1a_5$ 得 $(1+d)^2=1+4d$，即 $d=2$，故 $S_8=8a_1+28d=64$。
  \item[\textbf{13.}] $590$。$\mathrm C_{12}^5-\mathrm C_9^5-\mathrm C_8^5-\mathrm C_7^5+\mathrm C_5^5=792-126-56-21+1=590$。
  \item[\textbf{14.}] $5$。$AC=10$，$BC=10\sqrt3$，弦切角 $\angle BCD=\angle A=60^\circ$，故 $CD=5\sqrt3$，$BD=15$，再由 $DC^2=DE\cdot DB$ 得 $DE=5$。
  \item[\textbf{15.}] $16$。直线为 $x=4$，代入 $x=t^2$ 得 $t=\pm2$，交点 $(4,8)$，$(4,-8)$，故 $|AB|=16$。
  \item[\textbf{16.}] $a\leqslant8$。$|x-5|+|x+3|\geqslant8$ 恒成立，故不等式无解当且仅当 $a\leqslant8$。
\end{itemize}

\gkpart{解答题}
\begin{itemize}[leftmargin=2.6em]
  \item[\textbf{17.}] (1) $a=\dfrac12$。$f'(x)=2a(x-5)+\dfrac6x$，切线 $y=16a+(6-8a)(x-1)$ 过 $(0,6)$，得 $24a-6=6$；(2) 单调递增区间为 $(0,2)$ 与 $(3,+\infty)$，单调递减区间为 $(2,3)$；极大值 $f(2)=\dfrac92+6\ln2$，极小值 $f(3)=2+6\ln3$。由 $f'(x)=\dfrac{(x-2)(x-3)}x$ 判号即得
  \item[\textbf{18.}] (1) $\dfrac{18}{35}$。第一次恰摸 $1$ 个红球、第二次摸到白球，概率为 $\dfrac{\mathrm C_3^1\mathrm C_4^2}{\mathrm C_7^3}=\dfrac{18}{35}$；(2) $P(X=200)=\dfrac1{105}$，$P(X=50)=\dfrac2{105}$，$P(X=10)=\dfrac4{35}$，$P(X=0)=\dfrac67$，$E(X)=4$。三个奖项的概率分别为 $\dfrac1{35}\cdot\dfrac13$，$\dfrac1{35}\cdot\dfrac23$，$\dfrac{12}{35}\cdot\dfrac13$，期望为 $\dfrac{200+100+120}{105}=4$
  \item[\textbf{19.}] (1) $2\sqrt3$。由余弦定理 $AB=AD=2\sqrt3$，建系后 $\overrightarrow{AF}\cdot\overrightarrow{PB}=6-\dfrac{PA^2}2=0$，得 $PA=2\sqrt3$；(2) $\dfrac{3\sqrt7}8$。两半平面法向量可取 $(3,-3\sqrt3,-2\sqrt3)$ 与 $(-3,-3\sqrt3,2\sqrt3)$，$\cos\theta=\dfrac18$，故 $\sin\theta=\dfrac{3\sqrt7}8$
  \item[\textbf{20.}] (1) $C=\dfrac{3\pi}4$。由余弦定理 $\cos C=\dfrac{a^2+b^2-c^2}{2ab}=-\dfrac{\sqrt2}2$；(2) $\tan\alpha=1$ 或 $\tan\alpha=4$。展开条件得 $\tan^2\alpha-5\tan\alpha+4=0$，两根均满足 $\cos\alpha\neq0$
  \item[\textbf{21.}] (1) $\dfrac{x^2}{16}+\dfrac{y^2}8=1$。由 $e=\dfrac{\sqrt2}2$ 得 $b=c$，通径长 $\dfrac{2b^2}a=4$ 给出 $a=4$，$b^2=8$；(2) $\left(x-\dfrac{2\sqrt6}3\right)^2+y^2=\dfrac{16}3$ 或 $\left(x+\dfrac{2\sqrt6}3\right)^2+y^2=\dfrac{16}3$。由相切与 $PQ\perp P'Q$ 得 $x_0=2y_0$、$Q\left(\pm\dfrac{2\sqrt6}3,0\right)$，半径 $\dfrac{4\sqrt3}3$
  \item[\textbf{22.}] (1) $46$。$49$ 个 $\dfrac{m}{\sqrt k}$ 中只有 $\dfrac11=\dfrac2{\sqrt4}$、$\dfrac21=\dfrac4{\sqrt4}$、$\dfrac31=\dfrac6{\sqrt4}$ 三对重合，故共 $49-3=46$ 个；(2) $14$。$P_{15}$ 中 $1,3,6,10,15$ 首尾相接、相邻两数之和依次为 $4,9,16,25,16$，构成奇圈，无法二染色；$n=14$ 时可显式给出二分（见解析正文）
\end{itemize}
